\chapter{Fazit und Ausblick}
\label{chap:fazit}

\section{Zusammenfassende Bewertung}
\label{sec:bewertung}

Die Aufgabenstellung verlangte kein möglichst umfangreiches Produkt, sondern eine tragfähige und skalierbare Architektur. Gemessen an diesem Maßstab wurde das Ziel erreicht, denn das TeamBoard ist ein lauffähiges Microservice-System aus sieben Diensten und einem Protokoll-Adapter, das lokal mit einem Kommando startet und über eine Push-Pipeline in die Cloud ausrollt. Der Fokus lag hierbei vor allem auf den gewählten Architekturen und nicht auf der Implementierung möglicher Features, weswegen die Funktionalität  bewusst schlank blieb. Der in \vref{sec:funktionsumfang} zusammengefasste Funktionsumfang belegt beides, den umgesetzten Kern und die bewusst offen gehaltenen Lücken für zukünftige Weiterentwicklungen.

Die drei Leitfragen aus \vref{sec:kontext} bilden den roten Faden dieser Bewertung. An ihnen lässt sich am ehesten ablesen, wie gut die gewählte Architektur ihren eigenen Anspruch einlöst.

\paragraph{Verteilung von Echtzeit-Ereignissen.} Die erste Frage ist am deutlichsten beantwortet. Der ereignisgetriebene Kern aus \vref{sec:zusammenspiel} übernimmt die Verteilung von Echtzeit-Ereignissen von Anfang bis Ende. Ein fachliches Ereignis wird über die Outbox verlässlich auf den RabbitMQ-Topic-Exchange gelegt, von mehreren Diensten unabhängig konsumiert und über den Notification-Service bis in die offene WebSocket-Verbindung des Clients zugestellt. Dass ein zweites Browser-Fenster eine verschobene Aufgabe ohne Neuladen sieht, ist dabei kein Sonderweg der Oberfläche, sondern die sichtbare Wirkung dieses Kerns. Entscheidend ist dabei die Entkopplung, so dass ein zusätzlicher Empfänger, wie z.B. ein Webhook über den Plugin-Service, nur einen weiteren Konsumenten kostet und keine Änderung am Sender bedarf. Der Preis dieses Weges ist die Eventual Consistency nach außen, die allerdings über idempotente Verarbeitung und eine Dead-Letter-Queue beherrschbar bleibt.

\paragraph{Unabhängige Skalierbarkeit.} Die zweite Frage fällt zwiespältiger aus und verdient die ehrlichste Antwort. Die unabhängige Skalierbarkeit ist im Entwurf konsequent angelegt, im laufenden Betrieb aber nicht ausgespielt. Alle Dienste außer dem Notification-Service sind zustandslos, besitzen ihre eigene Datenbank und werden allein über Umgebungsvariablen konfiguriert, sodass sich jeder Dienst grundsätzlich einzeln vervielfachen ließe. Der bewusst gewählte Betrieb auf einer einzigen EC2-Instanz aus \vref{sec:betriebsmodell} macht von dieser Eigenschaft jedoch keinen Gebrauch. Diese Entscheidung war unter den Rahmenbedingungen des Projekts richtig, denn sie hält lokale und produktive Umgebung identisch und die Kosten minimal. Sie hat aber den erkennbaren Preis eines Single Point of Failure und eines Betriebs, in dem von der geforderten Skalierbarkeit nichts zu sehen ist. Der Zuschnitt der Dienste sorgt allerdings dafür, dass der Weg zu echter Skalierung ein Ausbau bestehender Kanten bleibt und keine Neuentwicklung erfordert. Insofern ist die Frage im Entwurf beantwortet, im Betrieb jedoch nur vorbereitet.

\paragraph{Konsistenter Datenzugriff.} Die dritte Frage beantwortet die Architektur über eine bewusste Zweiteilung. Wo eine sofortige verbindliche Antwort nötig ist, ermöglicht ein schmaler synchroner Kern die Konsistenz. Der Project-Service ist die alleinige Autorität für Berechtigungen, und der Board-Typ wird beim Anlegen eines Boards synchron aus der Board Registry aufgelöst. Überall sonst wird Konsistenz über Ereignisse hergestellt und ist damit letztlich konsistent statt sofort. Die Verlässlichkeit dieses Ansatzes ruht auf zwei Mechanismen: Zum Einen stellt das Outbox-Pattern sicher, dass ein Ereignis genau dann existiert, wenn die zugehörige Datenänderung committet wurde, und zum Anderen sorgen idempotente Konsumenten dafür, dass eine mehrfache Übermittlung keinen Schaden anrichtet. Die getrennten Datenbanken je Dienst ermöglichen diese Autonomie damit, dass fachübergreifende Sichten nicht per SQL-Join, sondern über Ereignisse und Schnittstellen zusammengeführt werden. Für den gewählten Zuschnitt ist dieser Tausch angemessen.

Über die drei Leitfragen hinaus hat sich der Grundzug des Entwurfs, Erweiterungen auf klar benannte Punkte zu beschränken, im Verlauf des Projekts mehrfach bestätigt. Am deutlichsten zeigt sich das an der nachträglichen Anbindung von Sprachmodell-Clients aus \vref{sub:mcp}. Der Ausbau des Auth-Service zum OAuth-2.1-Autorisierungsserver blieb auf den Auth-Service und den Protokoll-Adapter beschränkt, ohne dass ein Domain-Service davon berührt wurde oder ein Schema sich änderte. Ebenso ist die Laufzeit-Registrierung neuer Board-Typen aus \vref{sub:boardregistry} über die Nutzung des Notebooks aus \vref{sec:laufendes-system} im laufenden System nachvollziehbar. Solche Belege sind aus unserer Sicht der eigentliche Ertrag des Projekts, da sie den Anspruch der Architektur an konkreten Erweiterungen einlösen.

Zugleich bleiben die in \vref{sec:funktionsumfang} benannten Lücken bestehen und werden hier nicht schöngeredet. Zwei offene Punkte betreffen die Autorisierung, nämlich die noch nicht auf eine Administrationsrolle eingeschränkten Schreibrouten der Board Registry und die fehlende synchrone Mitgliedschaftsprüfung beim Abonnieren eines Board-Kanals. Beide sind bekannt, dokumentiert und im vorliegenden Stand bewusst nicht umgesetzt. Dass diese Lücken offen genannt werden, ist die Konsequenz aus einer während des Projekts wiederholt beobachteten Neigung der Dokumentation, dem Code vorauszulaufen, wie in \vref{sec:ki-programmerstellung} beschrieben.

\section{Ausblick}
\label{sec:ausblick}

Aus der Bewertung ergeben sich mehrere Richtungen für die Weiterentwicklung, die sich in drei Gruppen ordnen lassen. Sie bauen durchweg auf dem bestehenden Entwurf auf und verlangen keine grundlegende Neukonzeption.

Die erste und naheliegendste Gruppe schließt die bekannten Lücken und behebt bekannte Probleme. Die Schreibrouten der Board Registry wären auf eine Administrationsrolle einzuschränken, und der Notification-Service benötigt eine Zuordnung von Boards zu Projekten, um die Mitgliedschaft beim Abonnieren eines Kanals synchron prüfen zu können. Diese Zuordnung ließe sich über die bereits vorhandenen Ereignisse des Project-Service aufbauen und fügt sich damit in den ereignisgetriebenen Kern ein. Hinzu käme der noch fehlende Test des Tracing-Exportpfads, der einen eigenen Test-Collector erfordert. Diese Arbeiten erhöhen den bestehenden Funktionsumfang nicht, sondern bringen den realisierten Stand auf das im Entwurf angelegte Sicherheitsniveau.

Die zweite Gruppe spielt die Skalierbarkeit aus, die bislang nur vorbereitet ist. Der in \vref{sec:betriebsmodell} skizzierte Weg führt in drei Schritten weiter, die den Anwendungscode unberührt lassen. Zuerst würden die zustandsbehafteten Container in verwaltete Dienste ausgelagert, also PostgreSQL nach RDS, Redis nach ElastiCache, RabbitMQ nach Amazon MQ und der Objektspeicher nach S3. Anschließend zögen die Dienste selbst auf eine Container-Plattform wie ECS oder Fargate um, wobei die vorhandenen Produktions-Images unverändert weiterlaufen. Zuletzt verteilte ein Load Balancer mehrere Instanzen je Dienst, was seinen Nutzen erst dann entfaltet, wenn einzelne Dienste tatsächlich unterschiedlich stark belastet werden. Erst dieser letzte Schritt macht aus der entworfenen Skalierbarkeit eine im Betrieb sichtbare Eigenschaft und schließt damit die zweite Leitfrage auch auf der Betriebsseite.

Die dritte Gruppe betrifft die fachliche und die konzeptionelle Erweiterung entlang der bereits angelegten Erweiterungspunkte. Der Plugin-Service ist als offener Punkt für ausgehende Integrationen entworfen, sodass sich neben dem umgesetzten Webhook weitere Kanäle wie E-Mail oder Chat-Dienste als zusätzliche Plugins ergänzen ließen, ohne die Domain-Services zu berühren. Ebenso sind eingehende Webhooks noch offen. Bei der Darstellung der Boards ist mit der Presentation-Spezifikation und dem einen Erweiterungsfeld \texttt{presentation.view} der Rahmen für entfernt geladene Views bereits gesetzt. Der in ADR 0003 als Richtungsentscheidung beschriebene Renderer für einen Wert \texttt{view:\,\enquote{remote}} würde diesen Rahmen ausfüllen und echte Micro-Frontends erlauben, ohne das zugrunde liegende Modell zu ändern. Dass eine solche Erweiterung ohne Eingriff in die bestehenden Dienste möglich ist, war das Ziel des Entwurfs und bleibt die aussagekräftigste Prüfung seiner Tragfähigkeit.

Insgesamt zeigt der Ausblick, dass die tragenden Entscheidungen dieses Projekts nicht auf den erreichten Stand begrenzt sind, sondern einen Weg für die Weiterentwicklung vorzeichnen. Der bewusste Verzicht auf funktionale Breite zugunsten einer sauber geschnittenen Architektur zahlt sich genau an dieser Stelle aus, denn jede der genannten Erweiterungen ist ein Ausbau des Bestehenden und erfordert keine einschneidenden Anpassungen.
